\chapter{端点已统一，为什么整个区间仍未统一}\label{ch:intervals}
\section{现在已经知道每个元数的2k+1端点}
\begin{theorem}[所有元数的共同端点]
对任意整数 $k\ge2$，在 $n=2k+1$ 个位置上，有限 $k$ 平均可达性恰由 $G=1$ 判定。
\end{theorem}
因为 $\gcd(k,2k+1)=1$，一般的允许素因子条件在这里就是 $G=1$。奇数 $k$ 的证明来自前篇奇余量构造，其实际前提不需要 $k$ 为素数；偶数 $k$ 的补齐使用下面的返回。全部端点定理是无条件的。\src{S18}\src{S19}

这个结论给 $M(k)\le2k+1$，也给 $H(k)\le2k+1$。它不等于 $N(k)\le2k+1$，因为定义 $N$ 时还要求所有后续维数成功。

\section{偶数端点：先取得一个反射，再补除二}
写 $k=2b,n=4b+1$。在两个等大块和一个单点的核心上，已知两步返回的整数部分为
\[T_j=\begin{pmatrix}-k-1&j\\-n&2j\end{pmatrix},
\quad1\le j\le k,\qquad\det T_j=-j.
\]
实际矩阵是 $T_j/k$。保留单点的另一类两步返回为
\[F_h=\begin{pmatrix}1&(h-k)/(2k)\\0&h/k\end{pmatrix},
\qquad h\text{ 为偶数},\quad|h|\le k.
\]
两块等大，所以这次 $\Sigma=\left(\begin{smallmatrix}1&-1\\0&-1\end{smallmatrix}\right)$ 可以直接通过重新登记块角色实现。

令 $J=T_b\Sigma$，有
\[J^2=-bI,\qquad T_b^{-1}T_j=\Delta(j/b).
\]
若 $b$ 为偶数，取 $j=2b,b/2$ 就得到 $\Delta(2)$ 和 $\Delta(1/2)$。于是熟悉的交换子给上剪切，而
\[U(-1/2)T_1\Delta(-1/4)=-\tfrac12L(2n)
\]
给下剪切。二进同余控制器已经就绪。

若 $b\ge3$ 为奇数，先选 $\epsilon\in\{-1,1\}$ 使
\[a_0=(4b+1+\epsilon b)/4
\]
为整数：$b\equiv1\pmod4$ 取负号，$b\equiv3\pmod4$ 取正号。则 $1\le a_0\le k$，它模 $n$ 安全，而且
\[\operatorname{tr}(T_{a_0}F_{2\epsilon})=0.
\]
这个周期给 $F_2$ 的逆。再由
\[\Delta(1/b)F_2^{-1}=U((b-1)/2)
\]
和 $\Sigma$ 对上剪切的变号作用，同时取得 $\Delta(b)^{\pm1}$。初始深层随之建立在 $\Z[1/b]$ 上。

剩下的除二，与平方族采取同一思路。置
\[
J'=\Sigma J\Sigma=\begin{pmatrix}2b&-b\\n&-2b\end{pmatrix},
\quad H=\tfrac12\Delta(2)J'\Delta(2)
=\begin{pmatrix}b&-b\\n&-4b\end{pmatrix}.
\]
$H$ 和 $J'$ 的行列式都为 $b$，所以 $H(J')^{-1}$ 是旧环上的行列式一矩阵。用已有深层取逆，再由
\[J'\Delta(2)H^{-1}=2\Delta(1/2)
\]
取得除二。这样两个奇偶分支接到同一套 $U(\Z[1/2])$、$L(n\Z[1/2])$ 资源。$b=1$ 是前篇的二平均情形。

\section{端点为什么不需要GRH的小生成元定理}
从 $T_b^{-1}T_j$ 可以激活所有 $1\le j\le k$ 的安全倍率及其逆。模 $2k+1$ 的每个非零余数都有绝对值不超过 $k$ 的平衡代表；若余数为单位，这个代表也与模数互素。

因此这里的可用倍率已经覆盖全部单位群。平方族只供应到约 $\sqrt n$，需要小非剩余定理；端点供应到约 $n/2$，平衡代表就足够了。差别来自实际返回能够容纳的计数范围。

贝祖输送把核心送到 $(0^k,(-1)^k,k)$。先平均单点 $k$ 与 $k-1$ 个零，得到 $k$ 份一和一个剩余零；再两次各选 $k/2$ 份一和 $k/2$ 份负一，全体归零。

任意输入入口复用三角色见证转移。此处 $\gcd(k-1,2k+1)$ 至多含三，一个保留单点足够保护这个失效素数，所以偶数性不妨碍入口。

\section{奇合数的部分中间带也已经接通}
\begin{theorem}[奇元数的互素奇余量]
设 $k\ge3$ 为奇数，$1\le r<k$ 为奇数且 $\gcd(k,r)=1$。则 $n=2k+r$ 满足完整 $G=1$ 判据。
\end{theorem}
原素数证明在这部分实际使用的是 $k$ 模危险素数可逆、奇数配比和足够的选点容量。核对短区间筛选后，这些条件对奇合数也成立。入口得到
\[(x^k,(rz-x)^k,(-kz)^r).
\]
两个已知两步反射写成
\[J=\begin{pmatrix}0&d\\-1&0\end{pmatrix},\quad
Q=\begin{pmatrix}0&e\\-1&0\end{pmatrix},\quad
d=-r(k+r)/2,\quad e=d+n.
\]
$\gcd(k,r)=1$ 保证 $d,e$ 互素且均与 $n$ 互素，二者奇偶相反。它们给 $n$ 倍上、下剪切，再由前篇的第三方向与有限环消元取得所需同余层。

最后的小代表也仍可使用：一个单位的平衡代表 $u$ 若大于 $k$，以 $2u-\operatorname{sgn}(u)n$ 替换，绝对值变成小于 $r<k$；再除掉二的幂，得到可由真实数字返回实现的安全奇数。前篇的列准备与终端随后完成。

\begin{corollary}[奇素数幂的全部奇中间维数]
若 $k=p^a$，$p$ 为奇素数，则每个奇数 $n$ 满足 $2k<n<3k$ 时，完整条件都是 $\rad(G)\mid p$。
\end{corollary}
\begin{proof}
令 $r=n-2k$，$d=\gcd(k,r)$，$q=k/d$，$r_0=r/d$。商层有 $m=n/d=2q+r_0$ 个位置，$q,r_0$ 为互素奇数，因此上一条定理适用。

按同时好划分把原位置分成大小 $d$ 的块。因为 $d$ 是 $p$ 幂，商层仍只允许同一分母素数。商层每次选 $q$ 块，相当于原位置选 $qd=k$ 项平均，块均值精确模拟商操作。均值达到目标后，再以覆盖全部块的操作清除可能剩下的块内差异。
\end{proof}
例如九平均二十一、二十三、二十五元都已覆盖。这项结论没有包括偶数位置数。

\section{有公因子时，核心坐标必须先改正}
一般 $n=2k+r$，写 $d=\gcd(k,r)$、$k_0=k/d$、$r_0=r/d$。三块整数零和条件只保证载体值被 $k_0$ 整除，正确参数化为
\[(a^k,(r_0z-a)^k,(-k_0z)^r).
\]
三个块值的最大公因数是 $\gcd(a,z)$。本原时令 $y=a+k_0z$，两个独立差分为 $y$ 和 $2y-(n/d)z$，所以
\[G=\gcd(y,n/d).
\]
例如六平均、$r=3$ 时，核心 $(1^6,0^6,(-2)^3)$ 对应整数参数；若硬写载体为 $-6z$ 并要求 $z$ 为整数，就漏掉了它。\src{S15}

真正的危险模数只由 $n$ 中不整除 $k$ 的素数幂组成。$d$ 的素因子都允许出现在分母里，因此除去 $d$ 不改变危险部分。但在模整个 $n$ 下直接除以 $k$ 仍然不合法。

\section{共同入口不等于共同收尾}
对所有 $k\ge2,n\ge2k+1$，已经有移动见证、再合并的入口，得到两个 $k$ 重块与剩余 $r=n-2k$ 个单点。\src{S26}

它先保护使交换子失效的少数危险素数，再把其他见证通过CRT集中到一个角色，最后在保护位置外做平均。这解决了建立两个大块的问题。它并未自动把 $r$ 个单点压成一个等值块，也未证明之后总能继续降低核心维数。

若要得到 $N(k)\le2k+1$，还必须解决这些剩余参数与更多中间带。更进一步要得到 $N(k)=B(k)$，则还要填补最早成功点至 $2k+1$ 之间的范围。端点、部分奇中间带和大维数尾部是三块已建好的桥面，中间仍有未接通的区间。
